\documentclass[../main.tex]{subfiles}

\begin{document}
\chapter{Basic Concepts}

\section{Phase Space}
Consider a system living in a state space $S$, with time $\mathbb{R}$ as an evolution parameter. A \textbf{process} is a function $x: T \to S$ where $T \subseteq \mathbb{R}$ is a time interval. A process is called \textbf{deterministic} if we can determine the whole process if we know any point $x(t_0)$ for some $t_0 \in T$. In this context, we can call the state space $S$ the \textbf{phase space} of the system.

A process is called \textbf{finite-dimensional} if the phase space $S \subseteq \mathbb{R}^n$ for some $n \in \mathbb{N}$. A process is called \textbf{differentiable} if $S$ is a differentiable manifold.

\subsection{Integral Curves and Differential Equations}


\begin{plainblackenv}
  From now on we shall assume all objects encountered are smooth, i.e. continuously differentiable to all necessary orders.
\end{plainblackenv}
Often we knows a velocity vector field $v$ on the phase space $S$, and the whole problem of determining the process is to find an integral curve of the vector field. In coordinates, this results in ordinary differential equations (ODEs). Usually, as $S$ is a subset of $\mathbb{R}^n$, we usually identify any point $x \in S$ with its local standard coordinates $(x^1, \dots, x^n)$. But note that the whole space $S$ may not have a global coordinate because $S$ may not be diffeomorphic to a ball.

In the standard coordinates $(x^i)$ of $\mathbb{R}^n$, we can write $v = \sum_{i=1}^n v^i(x) \partial / \partial x^i$, and the integral curve $\gamma: T \to \mathbb{R}^n, t \mapsto (\gamma^i(t))$ satisfies the following system of ODEs:
\begin{equation}
  \frac{\mathrm{d} \gamma^i}{\mathrm{d} t} = v^i(\gamma(t)), \quad i = 1, \dots, n.
\end{equation}

If the vector field is time dependent, then we may rely on time-dependent vector fields theory or just promoting the time variable to a new coordinate, which results in an autonomous system of ODEs:
\begin{equation}
  \frac{\mathrm{d} \gamma^i}{\mathrm{d} t} = v^i(t, \gamma(t)), \quad i = 1, \dots, n.
\end{equation}
We can just interpret this as finding a new integral curve $\tilde{\gamma}: T \to \mathbb{R}^{n+1}, t \mapsto (t, \gamma(t))$ of the autonomous vector field $\tilde{v} = \partial / \partial t + \sum_{i=1}^n v^i(t, x) \partial / \partial x^i$.

\begin{definition}{Equilibrium Point}{Equilibrium Point}
  An \textbf{equilibrium point} of a vector field $v$ is a point $x_0$ such that $v(x_0) = 0$. In other words, the process starting at $x_0$ will stay at $x_0$ for all time.  
\end{definition}

At an equilibrium point $a$, the vector field $v$ vanishes, so a solution of $\gamma$ can be just $\gamma(t) = a$ for all $t$. In this case, the system is called \textbf{stationary}.

\begin{example}{One Dimensional Equations}{One Dimensional Equations}
  \begin{itemize}
    \item Consider the spacial uniform field $v = v(t) \partial / \partial x$ on $\mathbb{R}$. The integral curve is just
      \begin{equation*}
        \gamma(t) = \gamma(0) + \int_0^t v(\tau) \mathrm{d} \tau.
      \end{equation*}
    \item Consider the autonomous field $v = v(x) \partial / \partial x$ which is nonvanishing, then the integral curve is represented by
      \begin{equation}
        \int_{\gamma(0)}^{\gamma(t)} \frac{\mathrm{d} \gamma}{v(\gamma)} = t.
      \end{equation}
      As a nonvanishing smooth vector field on $\mathbb{R}$ is either positive or negative, both sides are well defined and increasing. So this determines $\gamma(t)$ implicitly as a function of $t$. It can be easily seen that this is the unique solution of the equation.
  \end{itemize}
\end{example}

\begin{remark}
  In analysis course, we have written informally that rearranging the infinite small would yield
  \begin{equation*}
    \frac{\mathrm{d} \gamma}{\mathrm{d} t} = v(\gamma) \implies \frac{\mathrm{d} \gamma}{v(\gamma)} = \mathrm{d} t \implies \int_{\gamma(0)}^{\gamma(t)} \frac{\mathrm{d} \gamma}{v(\gamma)} = \int_0^t \mathrm{d} t = t.
  \end{equation*}
  To be more formal, consider the integral curve $\gamma: T \rightarrow \mathbb{R}$, then we use the language of exterior derivative. We have the definition
  \begin{equation*}
    \gamma'(t) = \mathrm{d}_{\gamma(t)} \gamma \left( \frac{\partial }{\partial t} \bigg|_t \right) = v(\gamma(t)) \frac{\partial }{\partial x} \bigg|_{\gamma(t)}.
  \end{equation*}
  Therefore, we have
  \begin{equation*}
    (\gamma^* \mathrm{d} x)_t \left( \frac{\partial }{\partial t} \bigg|_t \right) = \mathrm{d}_{\gamma(t)} x \left( \mathrm{d}_{\gamma(t)} \gamma \left( \frac{\partial }{\partial t} \bigg|_t \right) \right) = \mathrm{d}_{\gamma(t)} x \left( v(\gamma(t)) \frac{\partial }{\partial x} \bigg|_{\gamma(t)} \right) = v(\gamma(t)).
  \end{equation*}
  Therefore, we have
  \begin{equation*}
    \gamma^* \mathrm{d} x = v \circ \gamma \mathrm{d} t, \qquad \gamma^* \left( \frac{\mathrm{d} x}{v(x)} \right) = \frac{1}{v \circ \gamma} \gamma^* \mathrm{d} x = \mathrm{d} t.
  \end{equation*}
  Integrating both sides, we have our result.
\end{remark}


\subsection{Examples}

\section{Vector Fields on the Line}
In this section we study a differential equation defined by a vector field on the line and equations with separable variables that reduce to such an equation.


\begin{theorem}{Existence and Uniqueness of Solutions for ODEs on the Line}{Existence and Uniqueness of Solutions for ODEs on the Line}
  Let $v$ be a smooth (continuously differentiable) function on some open interval $U \subseteq \mathbb{R}$, then the solution $\varphi$ of the equation $\dot{x} = v(x)$ with initial condition $(t_0, x_0)$:
  \begin{itemize}
    \item Exists for all $t_0 \in \mathbb{R}$ and $x_0 \in U$.
    \item Is unique, in the sense that in any two solutions coincides in some neighborhood of $t_0$.
  \end{itemize}
  The actual solution can be represented by the implicit formula
  \begin{equation}
    t - t_0 = \int_{x_0}^{\varphi(t)} \frac{\mathrm{d} \xi}{v(\xi)}, \qquad v(x_0) \neq 0.
  \end{equation}
  and $\varphi(t) = x_0$ for all $t$ if $v(x_0) = 0$.
\end{theorem}

\begin{remark}
  Consider the function $v(x) = x^{2 / 3}$ on $\mathbb{R}$, then the solutions $\varphi_1(t) = 0$ and $\varphi_2(t) = (t / 3)^3$ are two different solutions of the same initial value problem $\dot{x} = v(x), x(0) = 0$. This shows that the uniqueness part of the theorem fails if we drop the smoothness condition on $v$. In this case, a motion can reach the equilibrium point $0$ in finite time, and then we can choose to stay at $0$ or to leave $0$ after that.
\end{remark}

Now for the proof:
\begin{proof}
  Assume $\varphi$ is a solution. We have already construct the unique solution $\varphi$ when there are no equilibrium points. If there are equilibrium points, let $\varphi(t_0) = x_0$ be an equilibrium point, then we have a constant solution $\varphi(t) = x_0$ for all $t$. To see this is the unique solution, if there is $t_1$ such that $\varphi(t_1) 
\end{proof}

\end{document}
